% abstract.tex (Abstract)

\chapter*{Abstract}

Automatic modulation classification (AMC) is a key capability for intelligent receivers, spectrum monitoring, and cognitive radio networks where signals must be identified without prior coordination. Traditional AMC approaches rely on handcrafted features and expert rules, which can degrade under channel impairments and mismatched conditions. Recent advances in deep learning enable data-driven models to learn discriminative representations directly from raw I/Q samples, making AMC more robust in real-world environments. This report presents a baseline convolutional neural network (CNN) for classifying modulation types from the RadioML 2016.10a dataset, which contains labeled complex baseband signals across multiple modulation families and a wide range of signal-to-noise ratios (SNRs). The workflow includes per-sample power normalization, label encoding, and a stratified train/validation/test split. The CNN model uses stacked convolutional blocks with batch normalization, ReLU activations, max-pooling, and dropout to balance representation capacity and generalization. Training uses class-weighted cross-entropy and the AdamW optimizer with cosine learning rate scheduling. Performance is evaluated using overall accuracy, confusion matrices, and SNR-wise analysis. The baseline model achieves an overall test accuracy of approximately 0.574, with accuracy improving at higher SNR values and common confusions among closely related modulation families. The study confirms that compact CNNs are effective for AMC on raw I/Q data while highlighting the impact of SNR variability, class similarity, and dataset bias. The report concludes with a discussion of limitations and future directions, including data augmentation, deeper architectures, and domain adaptation for over-the-air generalization.
